% ltex: language=it
\documentclass{article}

% Shared preamble from project root
\usepackage{lecture}

\begin{document}

\begingroup
    \hypersetup{linkbordercolor=white, linkcolor=black} 
    \tableofcontents
    \newpage
\endgroup

\section{Descrizione del caso AUSL}\label{sec:descrizione_del_caso_ausl} % (fold)
Un'AUSL vuole valutare possibili soluzioni all'annoso problema dell'anagrafica assistita:
\begin{itemize}
    \item i residenti nella provincia di riferimento;
    \item i non residenti che per scelta o necessità usufruiscono di prestazioni;
\end{itemize}
AL momento i dati anagrafici degli assistiti sono mantenuti da più applicazioni, tra loro parzialmente incoerenti e ridondanti.
\begin{cimage}
    \includegraphics[width=.8\textwidth]{assets/ausl-initial.png}
\end{cimage}
\begin{description}
    \item[Scelta/Revoca del medico di base] È l'ufficio che si occupa di associare a tutti i residenti un medico di base/pediatra.
        La scelta/revoca deve essere fatta forzatamente presentandosi in ufficio.
        \item[CUP - Centro Unico Prenotazioni] È l'unico ufficio che gestisce le prenotazioni per tutte le prestazioni specialistiche.
            La prenotazione può essere fatta allo sportello oppure telefonicamente utilizzando il codice della ricetta del medico di base.
\end{description}
\begin{note}[Ogni struttura dispone di un ufficio accettazioni che regola ricoveri e dimissioni]
\end{note}
Il reparto di \textbf{pronto soccorso} ha un proprio sportello di accettazione vista l'urgenza dei ricoveri: i pazienti possono giungere al \textbf{PS} in stato confusionale o in stato di incoscienza e devono comunque essere registrati.
\emptyln{}
I \textbf{flussi regionali} permettono di comunicare con le altre AUSL le prestazioni svolte su assistiti non residenti:
\begin{itemize}
    \item in uscita sono spedite tutte le prestazioni/ricoveri effettuati da \textbf{NON} residenti;
    \item in entrata vengono ricevute tutte le prestazioni/ricoveri effettuati da residenti della provincia fatte in altre AUSL;
\end{itemize}
Come anagrafica principale è stata scelta quella dell'ufficio scelta/revoca del medico di base poiché tale scelta rappresenta un'operazione obbligatoria \textbf{per tutti e soli i residenti} e da fare di persona.
% section Descrizione del caso AUSL (end)
\end{document}
